\chapter{Consent Form (Possibly Remove This)}

\section*{Study background}
Large Language Models (LLMs) have become part of daily use amongst many during the last years. In order to ensure responsible use, education on capabilities and limitations of LLMs is necessary. This study aims to check the effectiveness of a mediation tool and how its use makes shifts in conceptual models.

\section*{Study Procedure}
The study is divided into 3 parts. The first part is an introductory interview, aimed at classifying your current understanding of how LLMs work, how you use them, and what consequences such use hold. The interview will last approximately 45 minutes.

The second part involves you getting access to the mediation tool, which you are free to use in your own time. In this part you will use the mediation tool to educate yourself on the concept of LLMs. You are free to use the tool for as long as you like.

The third part is a follow-up interview scheduled approximately 1 week after the introductory interview. This interview is aimed at getting your feedback on the mediation tool and views on LLMs. This interview will last approximately 30 minutes.

\section*{Voluntary Participation}
The participation in this study is completely voluntary, and you are free to stop the process and leave at any time and for any reason. There will be no negative consequences if you choose to withdraw your consent to participate. If you withdraw your consent, your data will not be used in any further research. You can request access to the data collected about you and you can also request it to be deleted. If you wish to withdraw consent from this study, please contact Alexander Höpner (contact details below).

\section*{Data Collection and Confidentiality}
The data collected and registered about you will only be used for the purpose of the project. An audio recording will be made of the interview; this will be used only to transcribe interview responses for the purpose of data analysis. Interview responses will be transcribed verbatim by the interviewer (researcher) onto a document, and any identifying information will be replaced with a participant number. Audio recording will be deleted after transcription is done.

All materials and results will be kept confidential; in particular, your name and any identifying or identified information will not be associated with the data. Your information will be stored securely and will be deleted in its entirety once this project is over. The results from this study will be publicly available and can be shared with other researchers. It will not be possible to identify you through the data that is shared with others. 

You have the right to access the information that is registered about you and to have any errors in this information corrected. You also have the right to the information about the data security measures that apply to the processing of the data. You can lodge a complaint about the processing of your data to the Norwegian Data Protection Authority and the institution’s Data Protection Officer.

\section*{Compensation}
For your participation, you are in the lottery to win a movie ticket.

\section*{Contact details}
This research is conducted by Alexander Höpner at the Department of Informatics, University of Bergen. The study is supervised by David Grellscheid. 

If you have any questions regarding this study, please contact Alexander Höpner at tir016@uib.no or +47 919 96 111.

You are free to contact the Data Protection Office at the University of Bergen (personvernombud@uib.no) if you have any questions about your rights as a research participant.